% Part: lambda-calculus
% Chapter: introduction
% Section: reduction

\documentclass[../../../include/open-logic-section]{subfiles}

\begin{document}

\olfileid{lam}{int}{red}
\olsection{ল্যাম্বডা পদের হ্রাস}

ল্যাম্বডা পদ নিয়ে কী করা যায়? তাকে সরল করা যায়। $M$ ও $N$ যেকোনো
ল্যাম্বডা পদ এবং $x$ যেকোনো চলরাশি হলে, $M$-এ~$x$-এর স্থানে $N$
প্রতিস্থাপন করে পাওয়া ফল নির্দেশ করতে আমরা $\Subst{M}{N}{x}$ ব্যবহার করতে
পারি। এটি চলরাশি-আবদ্ধতা এড়ানো প্রতিস্থাপন: প্রতিস্থাপনের পরে~$N$-এর মুক্ত
চলরাশিগুলিকে বদ্ধ করে ফেলত এমন $M$-এর বদ্ধ চলরাশিগুলির নাম আগে বদলে নিতে
হয়। যেমন,
\[
\Subst{(\lambd[w][xxw])}{yyz}{x} = \lambd[w][(yyz)(yyz)w].
\]

\begin{digress}
প্রতিস্থাপনের বিকল্প লিপিগুলি হলো $[N/x]M$, $[x/N]M$, এবং
$M[x/N]$-ও। সাবধান!
\end{digress}

স্বজ্ঞাতভাবে, $(\lambd[x][M])N$ এবং $\Subst{M}{N}{x}$-এর অর্থ একই;
প্রথম পদটির স্থানে দ্বিতীয়টি বসানোর ক্রিয়াকে \emph{$\beta$-সংকোচন}
বলে। $(\lambd[x][M])N$-কে একটি \emph{রিডেক্স} এবং
$\Subst{M}{N}{x}$-কে তার \emph{সংকুচিত রূপ} বলে। সাধারণভাবে, কোনো
উপপদের $\beta$-সংকোচনের সাহায্যে পদ~$P$-কে~$P'$-তে বদলানো সম্ভব হলে বলি,
$P$ \emph{এক ধাপে $P'$-তে $\beta$-হ্রাস পায়}, এবং লিখি $P \redone
P'$। $P$ থেকে কোনো সংখ্যক এক-ধাপের হ্রাসের সাহায্যে (একটিও না হতে পারে)
$P'$ পাওয়া গেলে, $P$ \emph{$P'$-তে $\beta$-হ্রাস পায়}; প্রতীকে,
$P \red P'$। যে পদকে আর $\beta$-হ্রাস করা যায় না, তাকে
\emph{$\beta$-অহ্রাস্য} বা \emph{$\beta$-স্বাভাবিক} বলে। প্রসঙ্গ
স্পষ্ট থাকলে আমরা ``$\beta$-হ্রাস পায়'' ইত্যাদির বদলে শুধু ``হ্রাস পায়''
বলব।

কয়েকটি উদাহরণ বিবেচনা করি।
\begin{enumerate}
\item আমাদের আছে
\begin{align*}
(\lambd[x][xxy]) \lambd[z][z] & \redone (\lambd[z][z])(\lambd[z][z]) y \\
& \redone (\lambd[z][z]) y \\
& \redone y.
\end{align*}
\item কোনো পদকে ``সরল'' করলে সেটি আরও জটিল হতে পারে:
\begin{align*}
(\lambd[x][xxy])(\lambd[x][xxy]) & \redone (\lambd[x][xxy])(\lambd[x][xxy])y \\
& \redone (\lambd[x][xxy])(\lambd[x][xxy])yy \\
& \redone \dots
\end{align*}
\item কোনো পদ অপরিবর্তিতও থেকে যেতে পারে:
\[
(\lambd[x][xx])(\lambd[x][xx]) \redone (\lambd[x][xx])(\lambd[x][xx]).
\]
\item আবার, কিছু পদকে একাধিক উপায়ে হ্রাস করা যায়; যেমন,
\[
(\lambd[x][(\lambd[y][yx]) z)] v \redone (\lambd[y][yv]) z
\]
সবচেয়ে বাইরের প্রয়োগটিকে সংকুচিত করে; এবং
\[
(\lambd[x][(\lambd[y][yx]) z)] v \redone (\lambd[x][zx]) v
\]
সবচেয়ে ভিতরেরটিকে সংকুচিত করে। তবে লক্ষ করো, এই ক্ষেত্রে উভয় পদই পরে
একই পদ~$zv$-তে হ্রাস পায়।
\end{enumerate}

শেষ উদাহরণের চূড়ান্ত ফলটি কাকতালীয় নয়; বরং এটি ল্যাম্বডা ক্যালকুলাসের
একটি গভীর ও গুরুত্বপূর্ণ ধর্মের দৃষ্টান্ত, যা ``চার্চ--রসার ধর্ম'' নামে
পরিচিত।

\end{document}
